\documentclass[11pt]{article}
\usepackage{mystyle}

\title{FIT1058 --- Chapter 5: Predicate Logic\\ \large Exercise Solutions}
\author{Alston}
\date{Semester 2, 2026}

\begin{document}
\maketitle


% ======================================================================
\begin{exercise}{5}
    (six degrees of separation)

    Consider the binary predicate $\text{knows}$ from p.\,68 in
    \S\,2.14. The domain of each of its arguments is the set of all
    people.

    It has been claimed that, in the human social network, the
    distance between any two people is at most 6. (See p.\,70 in
    \S\,2.15.)

    Write this claim in predicate logic, using just the predicate
    $\text{knows}$.
\end{exercise}

% ======================================================================
\begin{exercise}{7}
    For each of the following predicate logic statements, (i) identify
    the predicates, functions, variables (including a suggested domain
    for each) and constants used; (ii) state whether or not it is a
    proposition, and if it is, whether it is true or false.
    \begin{enumerate}[label=(\alph*)]
        \item $\forall x\ x^2\geq 0$
        \item $x^2<0$
        \item $\forall x\ 2x=x^2$
        \item $x<0 \Rightarrow \forall y\big((0<y)\Rightarrow(x<y)\big)$
        \item $\forall x\,\forall z\,\exists y\ (x<y)\wedge(y<z)$
        \item $\exists x\,\exists y\ \neg(x\Rightarrow y)$
    \end{enumerate}
\end{exercise}

% ======================================================================
\begin{exercise}{8}
    Consider the following predicate logic statement.
    \[
    \forall A\,\forall B\ (A\subseteq B)\Rightarrow(\mathcal{P}(A)\subseteq\mathcal{P}(B))
    \]
    What predicates, functions and variables does this statement use?
    Restate, in words, the theorem that is being stated here.
\end{exercise}

% ======================================================================
\begin{exercise}{9}
    Suppose that
    \begin{itemize}
        \item $W$ is a variable whose domain is the set of all English words,
        \item $L$ is a variable whose domain is the set of all English letters,
        \item $\text{WordContainsLetter}$ is a binary predicate which
            takes an English word as its first argument, an English
            letter as its second argument, and is True whenever that
            word contains that letter. For example,
            \[
            \text{WordContainsLetter}(\text{quizzical}, z) \text{ is True},
            \]
            \[
            \text{WordContainsLetter}(\text{quizzical}, e) \text{ is False}.
            \]
    \end{itemize}
    \begin{enumerate}[label=(\alph*)]
        \item Write the statement of Theorem 15 in predicate logic.
        \item Now suppose you also have a unary predicate $\text{Vowel}$
            whose domain is the set of English letters and which is
            True when its argument is a vowel. Now write the statement
            of Theorem 15 in predicate logic again, but using this new
            predicate to write it more compactly this time.
    \end{enumerate}
\end{exercise}

\begin{solution}
Write $\text{WordContainsLetter}(W,L)$ as $W(W,L)$.

\medskip
\noindent\textbf{(a)}
\[
\forall W\,\exists L\ \Big(\big(L=a\vee L=e\vee L=i\vee L=o\vee L=u\big)\wedge W(W,L)\Big)\ \vee\ W(W,y)
\]

\medskip
\noindent\textbf{(b)} Writing $\text{Vowel}(L)$ as $V(L)$:
\[
\forall W\,\exists L\ \Big(V(L)\wedge W(W,L)\Big)\ \vee\ W(W,y)
\]
or, with the disjunction brought inside the scope of $\exists L$:
\[
\forall W\ \Big(\exists L\ \big(V(L)\wedge W(W,L)\big)\ \vee\ W(W,y)\Big)
\]
\end{solution}

% ======================================================================
\begin{exercise}{10}
    Suppose you have the positive integer properties $\text{Even}$ and
    $\text{Prime}$, which are true precisely for even numbers and prime
    numbers respectively, and that you also have the binary predicate
    $\leq$ on $\mathbb{N}$. Write a predicate logic expression for the
    statement

    \begin{quote}
        There are infinitely many odd prime numbers but only finitely
        many even prime numbers.
    \end{quote}
\end{exercise}

\begin{solution}
Write $\text{Even}(n)$ as $E(n)$ and $\text{Prime}(n)$ as $P(n)$. By
the domain assumption, any element that is not even is odd, so
``odd'' is $\neg E(n)$, and ``odd prime'' is $P(n)\wedge\neg E(n)$.

\medskip
\noindent\textbf{Infinitely many odd primes.} There is no largest
odd prime -- for any bound $b$, some odd prime exceeds it:
\[
R = \forall b\,\exists n\ \big(P(n)\wedge\neg E(n)\wedge b\leq n\big)
\]

\medskip
\noindent\textbf{Only finitely many even primes.} There is a
largest even prime $m$, and every even prime is at most $m$:
\[
S = \exists m\ \Big(P(m)\wedge E(m)\ \wedge\ \forall n\big(P(n)\wedge E(n)\Rightarrow n\leq m\big)\Big)
\]

\medskip
\noindent\textbf{Final statement.}
\[
R\wedge S
\]
\end{solution}

% ======================================================================
\begin{exercise}{11}
    The ternary predicate $\text{Date}$ is defined for any
    $(d,m,y)\in\mathbb{N}\times\mathbb{N}\times\mathbb{N}$, and is True
    if $(d,m,y)$ represents a valid date in day-month-year format in
    the Gregorian calendar, and is False otherwise.

    Using $\text{Date}$ and $<$, write an expression in predicate logic
    which is True if and only if $(d_1,m_1,y_1)$ and $(d_2,m_2,y_2)$
    are both valid dates and the first date comes chronologically
    before the second.
\end{exercise}

\begin{solution}
Write $\text{Date}(d,m,y)$ as $D(d,m,y)$. $(d_1,m_1,y_1)$ comes
chronologically before $(d_2,m_2,y_2)$ exactly when
\[
y_1<y_2, \quad\text{or}\quad y_1=y_2 \wedge m_1<m_2, \quad\text{or}\quad y_1=y_2\wedge m_1=m_2\wedge d_1<d_2.
\]

So the answer is
\[
D(d_1,m_1,y_1)\wedge D(d_2,m_2,y_2)\wedge\Big((y_1<y_2)\vee\big((y_1=y_2)\wedge(m_1<m_2)\big)\vee\big((y_1=y_2)\wedge(m_1=m_2)\wedge(d_1<d_2)\big)\Big)
\]
\end{solution}

% ======================================================================
\begin{exercise}{12}
    Suppose you have
    \begin{itemize}
        \item the positive integer predicate $\leq$;
        \item the function $t\colon \{\text{programs}\}\times A^*\to\mathbb{N}$
            defined for all programs $P$ and input strings $x\in A^*$
            (where $A$ is an alphabet, \S\,1.5) by
            \[
            t(P,x) = \text{the time taken by program } P \text{ when run on input } x;
            \]
        \item the function $\text{len}\colon A^*\to\mathbb{N}_0$
            defined for any string $x\in A^*$ by
            \[
            \text{len}(x) = \text{the length of the string } x;
            \]
        \item binary functions for multiplication and exponentiation on
            the positive integers;
        \item a string variable $x\in A^*$;
        \item positive integer variables $c$ and $k$.
    \end{itemize}
    Write the following statement in predicate logic.
    \begin{quote}
        There exist integers $c$ and $k$ such that, for every input
        string $x$, the time taken by program $P$ on input $x$ is at
        most the product of $c$ and the $k$-th power of the length of
        $x$.
    \end{quote}
\end{exercise}

% ======================================================================
\begin{exercise}{13}
    There are many theorems that assert that there is a unique object
    that satisfies certain conditions. In this exercise, we look at how
    to make such statements using predicate logic.

    Suppose you want to say that there is a unique $x$ with property
    $P(x)$. It's not enough to just write
    \[
    \exists x\ P(x),
    \]
    because that only enforces existence without guaranteeing
    uniqueness. Sometimes, people use ``$\exists !$'' as a shorthand
    for ``there exists a unique''. With that shorthand,
    \[
    \exists ! x\ P(x)
    \]
    means

    \begin{quote}
        There exists a unique $x$ such that $P(x)$ holds.
    \end{quote}

    How would you write this statement just using the tools available
    to us in predicate logic, i.e., without using ``$!$''?
\end{exercise}

\begin{solution}
\[
\exists! x\ P(x) \;\equiv\; \exists x\ \Big(P(x)\wedge\big(\forall y\ (P(y)\Rightarrow x=y)\big)\Big)
\]
\end{solution}

% ======================================================================
\begin{exercise}{14}
    \begin{enumerate}[label=(\alph*)]
        \item Prove that
        \[
        \exists X\big(P(X)\wedge Q(X)\big) \text{ implies } \big(\exists X\ P(X)\big)\wedge\big(\exists X\ Q(X)\big)
        \]
        \item Prove that
        \[
        \forall X\big(P(X)\vee Q(X)\big) \text{ is implied by } \big(\forall X\ P(X)\big)\vee\big(\forall X\ Q(X)\big)
        \]
        in two ways: (i) reason it through directly, (ii) use the
        result from part (a) and what you've learned about the
        relationship between existential and universal quantifiers.
    \end{enumerate}
\end{exercise}


\end{document}